\documentclass[12pt,a4paper]{article}
\usepackage[utf8]{inputenc}
\usepackage[T1]{fontenc}
\usepackage[czech]{babel}
\usepackage{geometry}
\usepackage{titlesec}
\usepackage{enumitem}
\usepackage{xcolor}
\usepackage{mdframed}
\usepackage{booktabs}
\usepackage{array}
\usepackage{hyperref}
\usepackage{fancyhdr}
\usepackage{tcolorbox}

\geometry{margin=2.5cm}

% Barvy – dark mode
\definecolor{primary}{RGB}{220, 60, 60}
\definecolor{secondary}{RGB}{180, 80, 80}
\definecolor{highlight}{RGB}{40, 10, 10}
\definecolor{boxbg}{RGB}{35, 20, 20}
\definecolor{darkbg}{RGB}{18, 18, 18}
\definecolor{textmain}{RGB}{230, 220, 210}
\definecolor{textdim}{RGB}{160, 150, 140}

% Nadpisy
\titleformat{\section}{\large\bfseries\color{primary}}{}{0em}{}[\titlerule]
\titleformat{\subsection}{\normalsize\bfseries\color{primary}}{}{0em}{}

% Záhlaví/zápatí
\pagestyle{fancy}
\fancyhf{}
\rhead{\textcolor{primary}{\textit{Romeo a Julie – studijní poznámky}}}
\lhead{\textcolor{primary}{W. Shakespeare}}
\rfoot{\thepage}
\renewcommand{\headrulewidth}{0.4pt}

% Prostředí pro důležité informace
\tcbuselibrary{skins}
\newtcolorbox{dulez}{
  colback=boxbg,
  colframe=primary,
  coltitle=white,
  coltext=textmain,
  fonttitle=\bfseries,
  title=⚠ Důležité,
  arc=4pt
}
\newtcolorbox{citace}{
  colback=darkbg!80!black,
  colframe=secondary,
  coltext=textmain,
  fonttitle=\bfseries\itshape,
  arc=4pt,
  left=8pt
}
\newtcolorbox{tip}{
  colback=darkbg!80!black,
  colframe=blue!60!black,
  coltitle=white,
  coltext=textmain,
  fonttitle=\bfseries,
  title=💡 K zapamatování,
  arc=4pt
}

\begin{document}
\pagecolor{darkbg}
\color{textmain}

% Titulní strana
\begin{center}
  {\Huge \textbf{\textcolor{primary}{Romeo a Julie}}}\\[0.3em]
  {\large \textit{William Shakespeare}}\\[0.2em]
  {\normalsize \textcolor{textdim}{Studijní poznámky na test}}\\[1em]
  \rule{\linewidth}{1pt}
\end{center}

% ============================================================
\section{Autor – William Shakespeare}

\begin{itemize}[leftmargin=1.5em]
  \item \textbf{Narozen:} 23. dubna 1564, Stratford-upon-Avon, Anglie
  \item \textbf{Zemřel:} 23. dubna 1616, Stratford-upon-Avon
  \item \textbf{Období:} Renesance (anglická renesance, alžbětinské divadlo)
  \item \textbf{Hlavní působiště:} Londýn, divadlo \textbf{Globe Theatre}
  \item \textbf{Povolání:} dramatik, básník, herec, spolumajitel divadla
\end{itemize}

\subsection{Dílo Shakespeara (přehled)}

\begin{tabular}{ll}
  \toprule
  \textbf{Kategorie} & \textbf{Příklady} \\
  \midrule
  Tragédie & Hamlet, Othello, Macbeth, Král Lear, \textbf{Romeo a Julie} \\
  Komedie & Sen noci svatojánské, Zkrocení zlé ženy, Benátský kupec \\
  Historické hry & Richard III., Julius Caesar, Jindřich IV. \\
  Sonety & 154 sonetů (slavný sonet č. 18: „Mám tě přirovnat k letnímu dni?") \\
  \bottomrule
\end{tabular}

% ============================================================
\section{Základní informace o díle}

\begin{itemize}[leftmargin=1.5em]
  \item \textbf{Žánr:} \textbf{tragédie} (přesněji: romantická / milostná tragédie)
  \item \textbf{Vznik:} cca \textbf{1594–1596}
  \item \textbf{Původní název:} \textit{The Most Excellent and Lamentable Tragedy of Romeo and Juliet}
  \item \textbf{Předloha:} italská novela \textit{Luigi da Porta} (1530), poté zpracoval \textit{Matteo Bandello}; anglicky převedl \textit{Arthur Brooke} v básni \textit{The Tragicall Historye of Romeus and Juliet} (1562)
  \item \textbf{Místo děje:} \textbf{Verona} (Itálie), krátce Mantova
  \item \textbf{Čas děje:} neurčitě – renesanční Itálie
  \item \textbf{Forma:} drama psané převážně \textbf{blankversem} (nerýmovaný pětistopý jambický verš), částečně próza, prolog jako sonet
\end{itemize}

\begin{dulez}
  Romeo a Julie je považována za \textbf{nejslavnější milostnou tragédii} světové literatury.
  Ústřední téma: \textbf{láska vs. nenávist, osud, smrt jako důsledek rodové nenávisti}.
\end{dulez}

% ============================================================
\section{Postavy}

\subsection{Hlavní postavy}

\begin{tabular}{>{\bfseries}p{3.5cm} p{10cm}}
  \toprule
  Postava & Popis \\
  \midrule
  Romeo Montague & Syn Montaguů, mladík zamilovaný nejdříve do Rosalindy, pak do Julie. Vznětlivý, romantický, impulzivní. \\[0.3em]
  Julie Capuletová & Dcera Capuletů, 13 let, inteligentní, odhodlaná, v lásce statečnější než Romeo. \\[0.3em]
  Mercutio & Romeův přítel, vtipný a ironický, zemře v souboji s Tybaltem. \\[0.3em]
  Tybalt & Bratranec Julie, bojovný, nenávidí Montaguovy, zabije Mercutia. \\[0.3em]
  Friar Lawrence & Mnich, oddá Romea a Julii tajně, vymyslí záchranný plán s jedem. \\[0.3em]
  Chůva & Juliin důvěrník, pomáhá milencům se scházet, komická postava. \\[0.3em]
  Paris & Šlechtic, navrhovaný ženich Julie, zabit Romeem. \\[0.3em]
  Kníže Escales & Vládce Verony, snaží se zastavit boje rodin. \\[0.3em]
  Lord/Lady Capulet & Juliini rodiče, chtějí ji provdat za Parise. \\[0.3em]
  Lord/Lady Montague & Romeovi rodiče. \\
  \bottomrule
\end{tabular}

% ============================================================
\section{Děj – stručně po aktech}

\subsection{Prolog}
Sbor (Chorus) přednáší sonet – oznamuje, že oba milenci zahynou a teprve jejich smrt usmíří rodiny.

\begin{citace}
„Z krve dvou nepřátelských rodů stejně vznešených / zrodí se pár nešťastných milenců, / jejichž smrt pohřbí rodičovský hněv."
\end{citace}

\subsection{Akt I – Seznámení}
\begin{itemize}[leftmargin=1.5em]
  \item Ulice Verony: rvačka mezi sluhy Capuletů a Montaguů → kníže varuje pod trestem smrti.
  \item Romeo je smutný – zamilován do \textbf{Rosalindy}, která ho odmítá.
  \item Capuletové pořádají \textbf{maškarní ples}; Romeo se vloupá (aby viděl Rosalindu).
  \item Na plese se \textbf{Romeo setká s Julií} – okamžitá láska, nevědí, kdo druhý je.
  \item Zjistí, že jsou z nepřátelských rodin.
\end{itemize}

\subsection{Akt II – Milostná scéna, tajný sňatek}
\begin{itemize}[leftmargin=1.5em]
  \item \textbf{Balkonová scéna} – Romeo pod Juliiným balkónem, vyznání lásky.
  \item Julie: \textit{„Co je jméno? To, čemu říkáme růže, / vonět by stejně vonelo i jinak."}
  \item Romeo a Julie se \textbf{tajně vezmou} u Friara Lawrence.
  \item Chůva zprostředkovává komunikaci.
\end{itemize}

\subsection{Akt III – Tragický zvrat}
\begin{itemize}[leftmargin=1.5em]
  \item \textbf{Tybalt} vyvolá souboj s Romeem → Romeo odmítá (Tybalt je nyní jeho příbuzný).
  \item Mercutio bojuje za Romea → \textbf{Tybalt zabije Mercutia}.
  \item Romeo v hněvu \textbf{zabije Tybalta}.
  \item Kníže vyhostí Romea z Verony (\textbf{exil do Mantovy}).
  \item Rodiče chtějí provdat Julii za Parise → Julie zoufá.
\end{itemize}

\subsection{Akt IV – Plán s jedem}
\begin{itemize}[leftmargin=1.5em]
  \item Friar Lawrence dá Julii \textbf{uspávací nápoj} – bude vypadat mrtvá po 42 hodin.
  \item Plán: Romeo dostane zprávu, přijde do hrobky, Julie se probudí.
  \item Julie vypije nápoj v předvečer svatby s Parisem → nalezena „mrtvá".
  \item Pohřeb Julie v hrobce Capuletů.
\end{itemize}

\subsection{Akt V – Tragický konec}
\begin{itemize}[leftmargin=1.5em]
  \item Zpráva k Romeovi \textbf{nedorazí včas} (posel byl zadržen).
  \item Romeo slyší jen o Juliině smrti → koupí jed, spěchá do Verony.
  \item U hrobky zabije Parise, vejde dovnitř.
  \item Vidí „mrtvou" Julii → \textbf{vypije jed a zemře}.
  \item Julie se probudí, vidí mrtvého Romea → \textbf{probodne se jeho dýkou}.
  \item Rodiny Montaguů a Capuletů se nad mrtvými usmíří.
\end{itemize}

\begin{tip}
  Klíčový motiv: \textbf{osud a náhoda} – vše by bylo jinak, kdyby zpráva dorazila včas.
  Shakespeare záměrně ukazuje, jak malá náhoda může zničit životy.
\end{tip}

% ============================================================
\section{Témata a motivy}

\begin{tabular}{>{\bfseries}p{4cm} p{9cm}}
  \toprule
  Téma & Vysvětlení \\
  \midrule
  Láska & Čistá, absolutní láska překonávající vše – ale i sebezničující \\
  Nenávist a konflikt & Rodová nenávist jako zdroj tragédie \\
  Osud a náhoda & „Hvězdami prokletí milenci" – předurčenost záhuby \\
  Čas & Vše se odehrává velmi rychle (5 dní) – důraz na intenzitu \\
  Smrt & Smrt jako únik, ale i jako důsledek slepé nenávisti \\
  Rodina vs. jedinec & Konflikty povinnosti vůči rodině a osobní svobody \\
  Dospívání & Julie dospěje z dívky ve statečnou ženu \\
  \bottomrule
\end{tabular}

% ============================================================
\section{Literární forma a jazyk}

\begin{itemize}[leftmargin=1.5em]
  \item \textbf{Blankvers} – nerýmovaný jambický pentametr (základní forma)
  \item \textbf{Sonet} – prolog i scéna setkání Romea a Julie jsou ve formě sonetu
  \item \textbf{Próza} – komické postavy (chůva, sloužící)
  \item \textbf{Oxymóron} – hojně využíván: „milující nenávist", „chladný oheň", „nemocné zdraví"
  \item \textbf{Metafora, přirovnání} – Julie jako slunce, Romeo jako poutník
  \item \textbf{Dramatická ironie} – divák ví více než postavy (Julie není mrtvá)
\end{itemize}

\begin{citace}
\textbf{Slavné citáty:}\\[0.5em]
„Romeo, Romeo, proč jsi Romeo?" – Julie (II, 2)\\
„Co je jméno? Růže..." – Julie (II, 2)\\
„Pár hvězdami prokletých milenců." – Prolog\\
„Ó, jsem štěstěnou posměšný blázen!" – Romeo (III, 1)
\end{citace}

% ============================================================
\section{Historický a kulturní kontext}

\begin{itemize}[leftmargin=1.5em]
  \item \textbf{Renesance} – návrat k antickým ideálům, individualismus, humanismus
  \item \textbf{Alžbětinské divadlo} – Globe Theatre, jen mužští herci (i ženské role), žádné kulisy
  \item Dílo reflektuje \textbf{rodové spory} typické pro středověkou Itálii
  \item Téma \textbf{aranžovaných sňatků} – Julie byla zasnoubena Parisovi bez vlastní vůle
  \item \textbf{Vendetty} – italský zvyk krevní msty
\end{itemize}

% ============================================================
\section{Zpracování v kultuře}

\begin{tabular}{ll}
  \toprule
  \textbf{Dílo} & \textbf{Autor / rok} \\
  \midrule
  Opera \textit{Roméo et Juliette} & Charles Gounod, 1867 \\
  Balet \textit{Romeo a Julie} & Sergej Prokofjev, 1935 \\
  Muzikál \textit{West Side Story} & Leonard Bernstein / Stephen Sondheim, 1957 \\
  Film \textit{Romeo + Juliet} & Baz Luhrmann, 1996 (s L. DiCapriem) \\
  Film \textit{Romeo and Juliet} & Franco Zeffirelli, 1968 \\
  \bottomrule
\end{tabular}

% ============================================================
\section{Rychlá příprava na test – shrnutí}

\begin{tcolorbox}[colback=highlight, colframe=primary, coltitle=white, coltext=textmain, title=\textbf{Co si určitě zapamatovat}, fonttitle=\bfseries]
\begin{enumerate}[leftmargin=1.5em]
  \item \textbf{Autor:} William Shakespeare, anglická renesance, 1564–1616
  \item \textbf{Vznik:} cca 1594–1596, Globe Theatre v Londýně
  \item \textbf{Žánr:} romantická/milostná tragédie, drama, blankvers
  \item \textbf{Místo:} Verona (Itálie), krátce Mantova
  \item \textbf{Rodiny:} \textbf{Montague} (Romeo) vs. \textbf{Capulet} (Julie)
  \item \textbf{Klíčové postavy:} Romeo, Julie, Mercutio, Tybalt, Friar Lawrence, Chůva
  \item \textbf{Tragický zvrat:} Tybalt zabije Mercutia → Romeo zabije Tybalta → exil
  \item \textbf{Záchranný plán selže:} zpráva nedorazí → Romeo vypije jed, Julie se probodne
  \item \textbf{Morální poselství:} nenávist ničí i nevinné; rodiny se usmiřují až nad mrtvými
  \item \textbf{Hlavní témata:} láska, osud, nenávist, smrt, čas, konflikt rodiny a jedince
\end{enumerate}
\end{tcolorbox}

\vspace{1em}

\begin{tcolorbox}[colback=darkbg!80!black, colframe=blue!60!black, coltitle=white, coltext=textmain, title=\textbf{Časová osa – 5 dní děje}, fonttitle=\bfseries]
\begin{tabular}{>{\bfseries}cl}
  Neděle & Ples u Capuletů – setkání Romea a Julie \\
  Neděle (noc) & Balkonová scéna, slib lásky \\
  Pondělí & Tajný sňatek; smrt Mercutia a Tybalta; Romeo vyhoštěn \\
  Úterý & Romeo odjíždí do Mantovy; rodiče chtějí provdat Julii \\
  Středa & Julie vypije lektvar; pohřeb \\
  Čtvrtek (noc) & Romeo v hrobce → smrt obou milenců \\
\end{tabular}
\end{tcolorbox}

\vspace{2em}
\begin{center}
  \textcolor{textdim}{\small Poznámky sestaveny pro studijní účely · \textit{Romeo a Julie} · W. Shakespeare}
\end{center}

\end{document}
